\documentclass[12pt,letterpaper]{article}

% ── Packages ──────────────────────────────────────────────────────────────────
\usepackage[margin=1in]{geometry}
\usepackage{amsmath,amssymb,amsthm}
\usepackage{booktabs}
\usepackage{listings}
\usepackage{xcolor}
\usepackage{hyperref}
\usepackage{url}
\usepackage{graphicx}
\usepackage[expansion=false]{microtype}
\usepackage{setspace}
\usepackage{authblk}
\usepackage{abstract}
\usepackage{natbib}
\usepackage{fancyhdr}
\usepackage{titlesec}

% ── Theorem environments ───────────────────────────────────────────────────────
\newtheorem{theorem}{Theorem}
\newtheorem{definition}{Definition}
\newtheorem{corollary}{Corollary}

% ── Code listings ─────────────────────────────────────────────────────────────
\definecolor{codebg}{HTML}{0d1117}
\definecolor{codetext}{HTML}{e2e8f0}
\definecolor{codecomment}{HTML}{6b7280}
\definecolor{codekw}{HTML}{a855f7}
\definecolor{codestr}{HTML}{06b6d4}
\definecolor{codenumber}{HTML}{f97316}

\lstdefinestyle{sovereign}{
  backgroundcolor=\color{codebg},
  basicstyle=\ttfamily\footnotesize\color{codetext},
  commentstyle=\color{codecomment},
  keywordstyle=\color{codekw}\bfseries,
  stringstyle=\color{codestr},
  numberstyle=\tiny\color{codenumber},
  breakatwhitespace=false,
  breaklines=true,
  captionpos=b,
  frame=single,
  framesep=4pt,
  rulecolor=\color{codecomment},
  numbers=left,
  numbersep=8pt,
  showspaces=false,
  showstringspaces=false,
  showtabs=false,
  tabsize=2,
  xleftmargin=16pt,
}
\lstset{style=sovereign}

% ── Page style ─────────────────────────────────────────────────────────────────
\pagestyle{fancy}
\fancyhf{}
\rhead{\small Falsifiable Assurance via WORM Audit Chains}
\lhead{\small SnapKitty Sovereign AI Research}
\cfoot{\thepage}
\renewcommand{\headrulewidth}{0.4pt}

% ── Hyperref ──────────────────────────────────────────────────────────────────
\hypersetup{
  colorlinks=true,
  linkcolor=blue!70!black,
  citecolor=blue!70!black,
  urlcolor=blue!70!black,
  pdftitle={Falsifiable Assurance in Agentic AI Systems via Append-Only Cryptographic Audit Chains},
  pdfauthor={Ahmad Meta},
}

% ── Title ─────────────────────────────────────────────────────────────────────
\title{%
  \textbf{Falsifiable Assurance in Agentic AI Systems}\\
  \large via Append-Only Cryptographic Audit Chains
}

\author[1]{Ahmad Meta}
\affil[1]{SnapKitty Sovereign AI Research \\ \href{https://github.com/SNAPKITTYWEST}{github.com/SNAPKITTYWEST}}

\date{June 2026 \quad arXiv:cs.AI / cs.CR / cs.SE}

% ══════════════════════════════════════════════════════════════════════════════
\begin{document}
\maketitle
\thispagestyle{empty}

% ── Abstract ──────────────────────────────────────────────────────────────────
\begin{abstract}
Contemporary AI governance frameworks rely on \emph{claimed} behavioral
guarantees---audit logs that are mutable, access controls that can be overridden,
and compliance reports generated after the fact.
We propose a fundamentally different model: \textbf{falsifiable assurance},
in which every AI decision, tool invocation, and model output is immediately
sealed into a cryptographically chained, append-only ledger that cannot be
altered without detection.

We introduce the \textbf{WORM Audit Chain} (Write Once Read Many), a SHA-256
linked ledger applied at the boundary of each agent action in a multi-step
agentic workflow. Combined with the \textbf{Trust Deed}---a declarative
governance contract evaluated before every mutation---the system produces
behavioral guarantees that are \emph{externally verifiable} and
\emph{falsifiable by construction}.

We demonstrate the architecture on the \textbf{Frankenstein Workflow}, a
four-stage agentic pipeline (Brain $\rightarrow$ Hands $\rightarrow$ Legs
$\rightarrow$ Review) processing 1{,}247 real enterprise tasks across ERP,
code review, and CI/CD domains. We measure \emph{Assurance Density} (sealed
decisions per workflow step), \emph{Chain Integrity Rate} (proportion of seals
surviving tamper detection), and \emph{Trust Deed Rejection Rate} (governance
violations caught pre-execution).

Our results show that falsifiable assurance adds a median 4\,ms overhead per
agent action while providing a complete, tamper-evident behavioral record that
satisfies SOX~\S302, GDPR Article~5(2), and ISO~27001 Annex~A.12.4 audit
requirements \emph{without any post-hoc reporting step}.
The full implementation is released at \url{https://github.com/SNAPKITTYWEST}.
\end{abstract}

\vspace{0.5em}
\noindent\textbf{Keywords:} agentic AI, audit chains, cryptographic integrity,
falsifiable assurance, enterprise compliance, WORM ledger, governance

\newpage
\tableofcontents
\newpage

% ══════════════════════════════════════════════════════════════════════════════
\section{Introduction}
\label{sec:intro}

\subsection{The Auditability Gap}

Large language model deployments in enterprise settings face a fundamental
tension: the same flexibility that makes LLMs valuable---open-ended reasoning,
tool use, multi-step planning---makes them difficult to audit. When a model
decides to approve a purchase order, modify a financial record, or trigger a
downstream system, that decision is typically logged to a mutable database
that an administrator can alter, a log file that can be rotated, or not logged
at all.

We term the distance between what a system \emph{claims} it did and what can
be \emph{independently verified} the \textbf{auditability gap}.

Existing approaches to AI auditability fall into three categories:
\begin{enumerate}
  \item \textbf{Post-hoc explanation}
    (LIME~\citep{ribeiro2016lime}, SHAP~\citep{lundberg2017shap},
    attention visualization~\citep{jain2019attention})---explains model
    internals but not behavioral records.
  \item \textbf{Structured logging}---mutable; requires trust in the
    logging infrastructure itself.
  \item \textbf{Constitutional AI / RLHF}~\citep{bai2022constitutional,
    ouyang2022instructgpt}---shapes model behavior but produces no
    verifiable per-action record.
\end{enumerate}

None of these approaches produce what compliance frameworks actually require:
a tamper-evident record of \emph{every decision made, when it was made, and
what its inputs were}.

\subsection{Contributions}

This paper makes four contributions:
\begin{enumerate}
  \item \textbf{The WORM Audit Chain}: a SHA-256 linked, append-only ledger
    architecture for AI agent actions, with a formal tamper-detection proof
    (\S\ref{sec:worm}).
  \item \textbf{The Trust Deed}: a declarative pre-execution governance
    contract that blocks policy-violating actions before they execute
    (\S\ref{sec:trustdeed}).
  \item \textbf{Falsifiable Assurance}: a formal definition of verifiable AI
    behavioral guarantees with measurable metrics (\S\ref{sec:metrics}).
  \item \textbf{Empirical validation} on 1{,}247 real enterprise tasks, with
    mapping to SOX, GDPR, and ISO~27001 (\S\ref{sec:eval}).
\end{enumerate}

% ══════════════════════════════════════════════════════════════════════════════
\section{Background}
\label{sec:background}

\subsection{Agentic AI Workflows}

Multi-step agentic pipelines~\citep{yao2023react,schick2023toolformer,
chase2022langchain} decompose complex tasks into sequences of reasoning steps,
tool invocations, and state mutations. Each step is a potential audit point.
Existing frameworks treat these steps as ephemeral---computed, outputs passed
forward, intermediate state discarded.

Benchmarks such as AgentBench~\citep{liu2023agentbench} and
ToolBench~\citep{qin2023toolllm} evaluate agent \emph{capability} but not
agent \emph{accountability}.

\subsection{Cryptographic Audit Chains}

Append-only, hash-linked chains provide tamper-evidence: modifying any entry
invalidates all subsequent hashes, making tampering detectable by any party
holding the chain~\citep{nakamoto2008bitcoin,laurie2014certificatetransparency}.

We adapt this principle to per-action AI audit, eliminating the need for
distributed consensus. The chain is local, deterministic, and verifiable by
a third party given only the ledger file.

\subsection{Constitutional AI and Governance}

Constitutional AI~\citep{bai2022constitutional} and Sparrow~\citep{glaese2022sparrow}
define behavioral constraints at training time. We argue these are necessary
but insufficient: they shape priors but produce no per-execution record.
A well-aligned model that produces no verifiable record is unauditable;
falsifiable assurance provides the missing verification layer.

% ══════════════════════════════════════════════════════════════════════════════
\section{The WORM Audit Chain}
\label{sec:worm}

\subsection{Formal Definition}

Let $\mathcal{A} = (a_1, a_2, \ldots, a_n)$ be a sequence of agent actions.

\begin{definition}[WORM Seal]
A \emph{WORM seal} for action $a_i$ is defined as:
\[
  S_i = \mathrm{SHA256}\!\left(S_{i-1} \,\|\, t_i \,\|\, \mathrm{serialize}(a_i)\right)
\]
where $S_0 = 0^{256}$ is the genesis seal, $t_i$ is a Unix millisecond
timestamp, $\|$ denotes concatenation, and $\mathrm{serialize}$ is a
deterministic encoding of the action payload.
\end{definition}

The chain $\mathcal{C} = (S_1, S_2, \ldots, S_n)$ is tamper-evident:

\begin{theorem}[Tamper Detection]
\label{thm:tamper}
For any modified chain
$\mathcal{C}' = (S_1, \ldots, S_{k-1}, S'_k, \ldots, S'_n)$ where
$S'_k \neq S_k$, a verifier holding $\mathcal{C}$ and the original action
payloads can detect the modification at position $k$ in $O(n)$ time.
\end{theorem}

\begin{proof}
By induction. $S_k$ depends on $S_{k-1}$ and $a_k$ via SHA-256. Any
modification to $a_k$ produces $S'_k \neq S_k$ with overwhelming probability
under the collision-resistance assumption of SHA-256 (preimage resistance:
$2^{-256}$ per query under the random oracle model). Since $S_{k+1}$ is
computed from $S_k$, we have $S'_{k+1} \neq S_{k+1}$, and by induction
$S'_j \neq S_j$ for all $j \geq k$. A verifier replays the chain in $O(n)$
steps and detects the first mismatch at position $k$. \qed
\end{proof}

\begin{corollary}
Prepending or reordering entries is also detectable, since $S_i$ encodes
$S_{i-1}$: any change to the sequence order invalidates all subsequent seals.
\end{corollary}

\subsection{Implementation}

Algorithm~\ref{alg:worm} presents the seal procedure.

\begin{lstlisting}[language=Python, caption={WORM seal procedure}, label={alg:worm}]
def worm_seal(action, chain):
    entry = {
        "ts":        unix_ms(),
        "prev_hash": chain.last_hash,
        "payload":   serialize(action),
    }
    entry["this_hash"] = sha256(
        (entry["prev_hash"] + str(entry["ts"]) + entry["payload"]).encode()
    ).hexdigest()
    chain.append(entry)          # append-only: no update/delete
    chain.last_hash = entry["this_hash"]
    return entry
\end{lstlisting}

In our PostgreSQL~16 implementation, tamper resistance is enforced at the
schema level:

\begin{lstlisting}[language=SQL, caption={PostgreSQL WORM chain schema}]
CREATE TABLE worm_chain (
  id        BIGSERIAL PRIMARY KEY,
  ts        TIMESTAMPTZ NOT NULL DEFAULT now(),
  prev_hash CHAR(64)    NOT NULL,
  payload   JSONB       NOT NULL,
  this_hash CHAR(64) GENERATED ALWAYS AS (
    encode(sha256((prev_hash || payload::text)::bytea), 'hex')
  ) STORED
);

REVOKE UPDATE, DELETE ON worm_chain FROM PUBLIC;
REVOKE UPDATE, DELETE ON worm_chain FROM sovereign_app;
\end{lstlisting}

The \texttt{GENERATED ALWAYS AS} column makes the hash server-computed and
non-writable. The \texttt{REVOKE} statements eliminate the update/delete
attack surface at the database permission level, requiring privilege escalation
for any tampering attempt.

\subsection{Cross-System Chaining}

In multi-service pipelines (e.g., Node.js connector $\to$ Elixir IDE $\to$
Rust orchestrator), each service maintains its own chain. Cross-system integrity
is preserved by including the upstream service's latest seal hash in the first
payload of the downstream chain:

\[
  \texttt{node\_seal}_{42} \;\to\; \texttt{abzu\_seal}_{01}
    \;(\text{prev} = \texttt{node\_seal}_{42})
  \;\to\; \texttt{rust\_seal}_{07}
    \;(\text{prev} = \texttt{abzu\_seal}_{01})
\]

This creates a directed acyclic graph of evidence spanning the entire pipeline,
verifiable by replay.

% ══════════════════════════════════════════════════════════════════════════════
\section{The Trust Deed}
\label{sec:trustdeed}

\subsection{Pre-Execution Governance}

The Trust Deed is a declarative policy evaluated \emph{before} each state
mutation. Unlike post-hoc filtering, pre-execution governance prevents policy
violations from entering the WORM chain, maintaining the chain as a record of
\emph{permitted actions only}.

\begin{lstlisting}[language=Python, caption={Trust Deed evaluation}]
def trust_deed_check(action, rules):
    for rule in rules:
        if rule.condition(action):
            if rule.effect == BLOCK:
                worm_seal({type: "BLOCKED", action, rule: rule.id})
                raise GovernanceViolation(rule)
            elif rule.effect == WARN:
                worm_seal({type: "WARNED", action, rule: rule.id})
    return PERMITTED
\end{lstlisting}

\subsection{Rule Taxonomy}

We identify five categories of governance rules observed in enterprise AI
deployments (Table~\ref{tab:rules}).

\begin{table}[h]
\centering
\caption{Trust Deed rule taxonomy}
\label{tab:rules}
\begin{tabular}{@{}llll@{}}
\toprule
\textbf{Category} & \textbf{Example} & \textbf{Effect} & \textbf{Severity} \\
\midrule
Destructive Ops      & \texttt{rm -rf /}, \texttt{DROP TABLE}      & BLOCK          & CRITICAL \\
Scope Escalation     & Write outside authorized namespace          & BLOCK          & HIGH     \\
Financial Integrity  & Double-entry imbalance in journal entries   & BLOCK          & HIGH     \\
Elevated Privilege   & Production DB write from dev agent          & WARN+CONFIRM   & MEDIUM   \\
External Exfiltration& Unauthorized outbound API call              & BLOCK          & HIGH     \\
\bottomrule
\end{tabular}
\end{table}

\subsection{Double-Entry Enforcement}

For financial AI pipelines, accounting integrity is enforced as a Trust Deed
rule at the database level:

\begin{lstlisting}[language=SQL, caption={Double-entry constraint trigger}]
CREATE CONSTRAINT TRIGGER enforce_double_entry
  AFTER INSERT ON je_line
  DEFERRABLE INITIALLY DEFERRED
  FOR EACH ROW EXECUTE FUNCTION check_balance();
-- Rejects any commit where sum(debit) != sum(credit)
-- per journal_entry_id; rejection is WORM-sealed.
\end{lstlisting}

% ══════════════════════════════════════════════════════════════════════════════
\section{Falsifiable Assurance: Definition and Metrics}
\label{sec:metrics}

\begin{definition}[Falsifiable Assurance]
A behavioral guarantee $G$ about an agentic system $\mathcal{S}$ is
\emph{falsifiable} if and only if there exists a verification procedure $V$
such that, given the WORM chain $\mathcal{C}$ produced by $\mathcal{S}$,
$V(\mathcal{C})$ returns \textsc{true} if $G$ holds and \textsc{false} with
a specific counterexample if $G$ is violated.
\end{definition}

This contrasts with \emph{claimed assurance} (policy documents, training
objectives) which cannot be verified against a specific execution trace.

\subsection{Metrics}

\paragraph{Assurance Density (AD).}
\[
  \mathrm{AD} = \frac{|\mathcal{C}|}{|\mathcal{A}|}
\]
The ratio of WORM seals to agent actions. $\mathrm{AD} = 1.0$ indicates every
action is sealed. In our implementation, $\mathrm{AD} = 1.0$ by construction.

\paragraph{Chain Integrity Rate (CIR).}
\[
  \mathrm{CIR} = 1 - \frac{\text{tamper detections}}{\text{verification runs}}
\]
Measured by periodic chain replay. In 90 days of production operation, $\mathrm{CIR} = 1.0$.

\paragraph{Trust Deed Rejection Rate (TDRR).}
\[
  \mathrm{TDRR} = \frac{\text{blocked actions}}{\text{total attempted actions}}
\]
Measures governance effectiveness. In our ERP deployment, $\mathrm{TDRR} = 0.023$.

% ══════════════════════════════════════════════════════════════════════════════
\section{Experimental Evaluation}
\label{sec:eval}

\subsection{The Frankenstein Workflow}

We evaluate on the Frankenstein Workflow, a four-stage agentic pipeline:

\begin{itemize}
  \item \textbf{Brain}: Claude Sonnet 4.6 (AWS Bedrock)---structured
    architectural reasoning and task decomposition.
  \item \textbf{Hands}: FAISS + Neo4j + NetworkX---semantic retrieval from
    a sovereign corpus of 717K lines across 20+ languages.
  \item \textbf{Legs}: Granite Code 3B~\citep{mishra2024granite} running on
    a local NVIDIA RTX 3080---code generation and execution.
  \item \textbf{Review}: Claude Sonnet 4.6 (AWS Bedrock)---output
    verification against the Brain specification.
\end{itemize}

Each stage produces one or more WORM seals. The pipeline exposes a REST API
and receives events from a GitLab CI/CD webhook integration.

\subsection{Dataset}

We evaluate on 1{,}247 real enterprise tasks across three domains:
\begin{itemize}
  \item \textbf{ERP operations} (427 tasks): journal entries, purchase order
    matching, accounts receivable aging, inventory movements.
  \item \textbf{Code review} (512 tasks): push event analysis, merge request
    reviews, CI pipeline failure diagnosis.
  \item \textbf{Direct commands} (308 tasks): developer-issued
    \texttt{/snapkitty} commands in GitLab merge request comments.
\end{itemize}

\subsection{Results}

Table~\ref{tab:results} presents quantitative results.

\begin{table}[h]
\centering
\caption{Evaluation results on 1{,}247 enterprise tasks}
\label{tab:results}
\begin{tabular}{@{}lr@{}}
\toprule
\textbf{Metric} & \textbf{Value} \\
\midrule
Total WORM seals generated          & 9{,}847 \\
Assurance Density (AD)              & 1.00 \\
Chain Integrity Rate (CIR)          & 1.00 \\
Trust Deed Rejection Rate (TDRR)    & 0.023 \\
\midrule
Median seal latency                 & 4\,ms \\
p99 seal latency                    & 11\,ms \\
Brain median latency                & 14{,}200\,ms \\
Legs median latency                 & 1{,}077\,ms \\
End-to-end median latency           & 27{,}800\,ms \\
\midrule
Double-entry violations caught (ERP) & 31 \\
Destructive command attempts blocked & 7 \\
\bottomrule
\end{tabular}
\end{table}

The 4\,ms median seal latency represents a 0.014\% overhead on the 27.8\,s
end-to-end workflow---negligible in practice.

\subsection{Compliance Mapping}

Table~\ref{tab:compliance} maps WORM chain properties to enterprise regulatory
requirements.

\begin{table}[h]
\centering
\caption{Regulatory compliance mapping}
\label{tab:compliance}
\begin{tabular}{@{}p{2.8cm}p{4.5cm}p{5cm}@{}}
\toprule
\textbf{Regulation} & \textbf{Requirement} & \textbf{Satisfied By} \\
\midrule
SOX \S302
  & Officer certification of financial controls
  & WORM chain provides independent verification basis \\[4pt]
SOX \S404
  & Management assessment of internal controls
  & Trust Deed + TDRR = falsifiable control evidence \\[4pt]
GDPR Art.\ 5(2)
  & Accountability: demonstrate compliance
  & Chain replay produces complete processing record \\[4pt]
GDPR Art.\ 22
  & Automated decision-making documentation
  & Every AI decision sealed with inputs and outputs \\[4pt]
ISO 27001 A.12.4
  & Event logging integrity
  & SHA-256 chained, append-only, cryptographically sealed \\
\bottomrule
\end{tabular}
\end{table}

% ══════════════════════════════════════════════════════════════════════════════
\section{Discussion}
\label{sec:discussion}

\subsection{Limitations}

The WORM chain provides tamper-evidence, not tamper-prevention at the hardware
level. A sufficiently privileged attacker (database administrator, system
administrator) could delete the chain file. We mitigate this through:
\begin{itemize}
  \item Off-site chain replication (S3-compatible object storage, sealed per sync)
  \item Cross-system hash anchoring (upstream seal embedded in downstream chain)
  \item Periodic third-party verification via the public verification API
\end{itemize}

The Trust Deed enforces declared policies but cannot catch undeclared threat
models. Novel attack patterns not anticipated in the governance rule set will
not be blocked.

\subsection{Relation to AI Safety}

Falsifiable assurance is a \emph{complement} to, not a replacement for,
alignment research. Constitutional AI and RLHF shape model priors; falsifiable
assurance creates a post-hoc verification layer. Both are necessary: a
well-aligned model that produces no verifiable record is unauditable; a fully
audited misaligned model is transparent but still harmful.

\subsection{Enterprise Deployment}

The 4\,ms median seal latency represents negligible overhead on workflows with
14--28 second end-to-end latency. We observe no throughput degradation at
50 concurrent workflow requests.

% ══════════════════════════════════════════════════════════════════════════════
\section{Conclusion}
\label{sec:conclusion}

We presented \emph{falsifiable assurance}---a formal framework for producing
cryptographically verifiable behavioral records from agentic AI systems.
The WORM Audit Chain and Trust Deed together provide:

\begin{enumerate}
  \item A tamper-evident record of every AI decision (Theorem~\ref{thm:tamper})
  \item Pre-execution governance that prevents policy violations
  \item Measurable, auditable metrics: AD, CIR, TDRR
  \item Direct mapping to SOX, GDPR, and ISO~27001 requirements
\end{enumerate}

The architecture is language-agnostic (implemented in Rust, Node.js, Elixir,
and PostgreSQL), adds 4\,ms median latency, and requires no changes to the
underlying model or training procedure.

We release the full implementation under the Sovereign Source License at
\url{https://github.com/SNAPKITTYWEST}.

% ── References ────────────────────────────────────────────────────────────────
\newpage
\bibliographystyle{abbrvnat}
\bibliography{paper}

% ══════════════════════════════════════════════════════════════════════════════
\appendix

\section{WORM Chain Verification Algorithm}
\label{app:verify}

\begin{lstlisting}[language=Python, caption={Chain verification procedure}]
import json, hashlib

def verify_chain(chain_file):
    entries = [json.loads(l) for l in open(chain_file)]
    prev = '0' * 64
    for i, entry in enumerate(entries):
        src = (prev + str(entry['ts']) + entry['payload']).encode()
        expected = hashlib.sha256(src).hexdigest()
        if entry['this_hash'] != expected:
            return False, f"Tamper detected at entry {i}: " \
                          f"expected {expected[:16]}..., " \
                          f"got {entry['this_hash'][:16]}..."
        prev = entry['this_hash']
    return True, f"Chain intact: {len(entries)} entries verified"
\end{lstlisting}

\section{Trust Deed Rule DSL}
\label{app:trustdeed}

\begin{lstlisting}[language={},caption={Trust Deed governance rule format (JSON)}]
{
  "version": "1.0",
  "rules": [
    {
      "id":          "no-worm-delete",
      "description": "Prevent deletion of WORM chain entries",
      "pattern":     "TRUNCATE worm_chain|DELETE FROM worm_chain",
      "effect":      "BLOCK",
      "severity":    "CRITICAL"
    },
    {
      "id":          "double-entry",
      "description": "All journal entries must balance",
      "check":       "sum(debit) == sum(credit) per journal_entry_id",
      "effect":      "BLOCK",
      "severity":    "HIGH"
    },
    {
      "id":          "no-root-delete",
      "description": "No recursive deletion of root filesystem",
      "pattern":     "rm -rf /|format c:|del /s \\*",
      "effect":      "BLOCK",
      "severity":    "CRITICAL"
    }
  ]
}
\end{lstlisting}

\section{Multi-Service Chain Graph}
\label{app:graph}

In a three-service deployment, the cross-system WORM graph takes the following
structure:

\begin{lstlisting}[caption={Cross-system chain anchoring}]
# Service A (Node.js) — seal 42
{prev: "0"*64, ts: 1719532800000, payload: {...}, this_hash: "a1b2c3..."}

# Service B (Elixir ABZU IDE) — seals anchor to A's last hash
{prev: "a1b2c3...", ts: 1719532801000, payload: {...}, this_hash: "d4e5f6..."}

# Service C (Rust orchestrator) — anchors to B
{prev: "d4e5f6...", ts: 1719532802000, payload: {...}, this_hash: "g7h8i9..."}

# Verifier replays A -> B -> C in sequence.
# Any gap in prev_hash linkage indicates a missing or tampered record.
\end{lstlisting}

\end{document}
